\documentclass{phys400proposal}
\usepackage[dvipsnames]{xcolor}
\proposaltitle{Automated Construction of Multi-Element MEAM Interatomic Potentials via DFT-Informed Neural Network Optimization}
\studentname{Can Durmaz}
\studentid{2587772}
\advisor{Prof. Dr. Hande Toffoli}
\studyfield{Computational Materials Science}
\studytype{Computational}
\submissiondate{\today}

\begin{document}

\makeproposalcover

%\tableofcontents
\newpage

\begin{boxedabstract}
\small
Molecular dynamics (MD) simulations of multi-component metallic alloys require interatomic potentials that accurately capture the interactions between all constituent elements. The Modified Embedded Atom Method (MEAM) provides a well-established framework for describing metallic bonding, yet published MEAM parameter sets typically cover only small subsets of elements, and no single file spans the full composition space of industrially relevant alloys such as Al~7075 (Al--Zn--Mg--Cu). When elements do not share a common MEAM parameterization, cross-interaction terms are unavailable and simulations cannot proceed. This project develops an automated pipeline that bridges this gap by combining first-principles density functional theory (DFT) calculations with neural-network surrogate optimization. Quantum Espresso is used to compute reference properties---equilibrium lattice constants, cohesive energies, bulk moduli, elastic constants, and binary formation energies---for each element and element pair. These quantities initialize MEAM parameters through analytic mapping relations derived from the Rose universal equation of state. A neural network surrogate is then trained on LAMMPS-evaluated elastic properties (Young's modulus and Poisson's ratio) across randomized alloy compositions, enabling gradient-based inverse optimization of the MEAM parameter set to reproduce target mechanical behavior. The resulting pipeline accepts an arbitrary list of elements, automatically discovers and merges existing MEAM libraries, fills missing cross-terms from DFT, and produces optimized multi-element potential files ready for large-scale MD simulations.

\vspace{0.5em}
\end{boxedabstract}

\introduction
Atomistic simulations of metallic alloys rely on interatomic potentials that balance computational efficiency with physical accuracy. Among classical potential formalisms, the Modified Embedded Atom Method \citep{baskes1992} extends the original EAM by incorporating angular dependence of bonding, making it suitable for a broad range of metallic and covalent systems. MEAM potentials have been widely adopted for studying mechanical properties, defect energetics, and phase stability in metals and alloys \citep{lee2001}.

However, a persistent challenge in multi-component alloy simulation is the availability of consistent MEAM parameterizations. Published potential files typically cover small element subsets---for example, a ternary Al--Mg--Zn set or a quinary Al--Si--Cu--Mg--Fe set---and the cross-interaction parameters between elements from different files are undefined. For an alloy such as Al~7075, which contains Al, Zn, Mg, Cu, and trace elements (Cr, Mn, Fe, Si, Ti), no single published MEAM file spans the full composition. Naively combining parameters from separate sources without proper cross-terms leads to unphysical results.

Traditional approaches to this problem involve laborious manual fitting of MEAM parameters against experimental or \textit{ab initio} databases \citep{jelinek2012}. More recently, machine-learning interatomic potentials (MLIPs) trained on DFT data have shown remarkable accuracy \citep{behler2007, shapeev2016}, but they require large training datasets and lack the interpretable functional form of MEAM. Neural-network interatomic potentials (NNIPs) represent an intermediate approach that can either replace or augment classical potentials \citep{mishin2021}.

This project takes a hybrid approach: rather than replacing MEAM with a purely data-driven potential, we use neural networks as surrogate models to optimize MEAM parameters so that the resulting classical potential reproduces DFT-quality mechanical properties. This preserves the computational efficiency and transferability of the MEAM formalism while leveraging DFT accuracy for parameter initialization and neural-network optimization for automated fitting.

\begin{table}[htbp]
	\centering
	\caption{Available MEAM potential libraries and their element coverage. No single file covers all elements needed for Al~7075 simulation.}
	\label{tab:meam_coverage}
	\begin{tabular}{|l|l|c|}
		\hline
		\textbf{Library File} & \textbf{Elements} & \textbf{Count} \\
		\hline
		\texttt{library\_AlMgZn.meam} & Al, Mg, Zn & 3 \\
		\texttt{library\_AlSiCuMgFE.meam} & Al, Si, Cu, Mg, Fe & 5 \\
		\texttt{library\_CuMo.meam} & Cu, Mo & 2 \\
		\texttt{library\_FeMnNiTiCuCrCoAl.meam} & Fe, Mn, Ni, Ti, Cu, Cr, Co, Al & 8 \\
		\hline
	\end{tabular}
\end{table}

As shown in Table~\ref{tab:meam_coverage}, the existing MEAM libraries in this project cover 12 elements in total, but each file spans only a subset. Elements such as Al and Cu appear in multiple files with potentially different parameterizations, while cross-terms between elements from different files (e.g., Zn--Fe or Mg--Cr) are entirely absent.

\researchquestion
Can an automated pipeline combining density functional theory reference calculations with neural-network surrogate optimization produce multi-element MEAM interatomic potentials that accurately reproduce the elastic properties of arbitrary metallic alloy compositions, even when the constituent elements do not share a common published MEAM parameterization?

\methodology
The pipeline consists of five stages, each implemented as a modular Python component orchestrated by a central driver script.

\textbf{Stage 1: Configuration Generation.}
Random alloy compositions are generated by sampling element fractions at $10^{-4}$ precision from the set of 12 available elements (Al, Co, Cr, Cu, Fe, Mg, Mn, Mo, Ni, Si, Ti, Zn). Each configuration specifies composition fractions, simulation box size, temperature, and integration parameters, and is stored as a JSON file.

\textbf{Stage 2: Molecular Dynamics Evaluation.}
Each composition is simulated in LAMMPS \citep{thompson2022} using the MEAM potential formalism. The elastic tensor is computed via a strain--stress protocol: after energy minimization under anisotropic barostat conditions, uniaxial strains of $\delta = 0.001$ are applied along three orthogonal directions and the resulting stress tensors are used to extract the Voigt-averaged elastic constants $C_{11}$ and $C_{12}$ via central differences. Young's modulus $E$ and Poisson's ratio $\nu$ are then computed as:
\begin{equation}
	E = \frac{(C_{11} - C_{12})(C_{11} + 2C_{12})}{C_{11} + C_{12}}, \qquad \nu = \frac{C_{12}}{C_{11} + C_{12}}
	\label{eq:elastic}
\end{equation}
Configurations yielding $\nu < 0$ are discarded as unphysical.

\textbf{Stage 3: DFT Reference Calculations.}
For each element, Quantum Espresso \citep{giannozzi2009} is used via the Atomic Simulation Environment (ASE) to compute: (i) the equation of state from 7 strained unit cells fit to a Birch--Murnaghan equation, yielding the equilibrium lattice constant $a_0$, cohesive energy $E_\mathrm{coh}$, and bulk modulus $B$; (ii) elastic constants $C_{11}$ and $C_{12}$ from stress--strain calculations; and (iii) isolated atom energies for true cohesive energy determination. For each binary pair, formation energies are computed in reference crystal structures (L1$_2$ for FCC, B2 for BCC). Spin-polarized calculations (nspin = 2) are employed for magnetic elements (Fe, Cr, Mn). PAW-PBE pseudopotentials are used with plane-wave cutoffs of 40~Ry (wavefunction) and 320~Ry (charge density).

\textbf{Stage 4: MEAM Parameter Initialization and Merging.}
DFT-computed properties are mapped onto MEAM parameters using established relations: the sublimation energy $E_\mathrm{sub} \leftarrow E_\mathrm{coh}$, the equilibrium lattice parameter $a_\mathrm{lat} \leftarrow a_0$, and the Rose equation shape parameter $\alpha = \sqrt{9B\Omega / E_\mathrm{coh}}$, where $\Omega$ is the atomic volume. Cross-interaction parameters $E_c(i,j)$ are derived from binary formation energies. Existing MEAM library files are automatically discovered, parsed, and merged; missing parameters are filled from DFT results.

\textbf{Stage 5: Neural Network Surrogate Optimization.}
The DFT-initialized MEAM parameters are refined using a neural network surrogate model. Random perturbations of the parameter set are sampled, and each sample is evaluated in LAMMPS across all target compositions. A feedforward neural network (architecture: 20--20--10 hidden units) is trained to map MEAM parameters to predicted elastic properties ($E$, $\nu$). Gradient-based inverse optimization through the trained surrogate identifies the parameter set that minimizes the discrepancy between predicted and target elastic properties. The optimized parameters are validated by direct LAMMPS evaluation, with a pass criterion of average error below 10\%.

\expectedresults
The primary expected outcome is a set of optimized multi-element MEAM potential files that reproduce the elastic properties (Young's modulus and Poisson's ratio) of metallic alloy compositions with an average error below 10\% relative to MD-computed targets initialized from DFT data. Specifically:

\begin{enumerate}
	\item \textbf{Unified MEAM potentials:} Complete library and parameter files covering all 12 elements (Al, Co, Cr, Cu, Fe, Mg, Mn, Mo, Ni, Si, Ti, Zn) with consistent cross-interaction terms, enabling simulations of arbitrary alloy compositions within this element set.

	\item \textbf{DFT-quality initialization:} Per-element properties (lattice constants, cohesive energies, bulk moduli, elastic constants) and binary formation energies that agree with published DFT and experimental values, providing a physically grounded starting point for optimization.

	\item \textbf{Neural network surrogate accuracy:} A trained surrogate model that predicts the elastic response of MEAM parameter sets with sufficient fidelity to guide inverse optimization, reducing the number of expensive LAMMPS evaluations by an order of magnitude compared to brute-force parameter sweeps.

	\item \textbf{Validation on Al~7075 compositions:} Demonstration that the optimized potentials correctly predict the elastic behavior of Al~7075-type compositions (Al--Zn--Mg--Cu with trace Cr, Mn, Fe, Si, Ti), a system that cannot be simulated with any single existing MEAM file.
\end{enumerate}

\begin{semestergantt}
\workpackage{Literature Review}{1}{3}
\workpackage{DFT Reference Calculations}{2}{5}
\workpackage{MEAM Initialization Pipeline}{4}{7}
\workpackage{NN Surrogate Training}{6}{9}
\workpackage{Inverse Optimization}{8}{11}
\workpackage{Validation \& Testing}{10}{13}
\workpackage{Writing \& Documentation}{11}{15}
\workpackage{Poster Preparation}{13}{15}
\end{semestergantt}

\referencessection

\end{document}
